\section{Introduction}
\label{sec:intro}

Deep reinforcement-learning agents are typically benchmarked on
short-horizon tasks with rich, dense reward signals---most commonly the
Atari Learning Environment~\citep{bellemare2013arcade,mnih2015atari}
and procedurally-generated grid worlds~\citep{cobbe2020procgen,
matthews2024craftax,samvelyan2021minihack}.  These environments stress
exploration, sample efficiency, and generalisation, but few of them
require an agent to maintain coherent intent across the millions of
steps that real long-horizon games---role-playing games, strategy
games, or open-world games---demand.

\emph{Pok\'emon Red}, a 1996 Game Boy role-playing game, sits at an
unusually informative point in the difficulty landscape.  Its
state-action space is small enough to render with a Game Boy emulator
at thousands of frames per second on consumer hardware, but the game
contains long causal chains (deliver Oak's parcel before exiting Viridian;
beat Brock with a fire- or water-resistant team; navigate Mt. Moon
without losing the run) that punish agents which lack temporal
abstraction.  The 18 event flags between the player's bedroom and the
Boulder Badge encode the kind of structured progress that the dense
reward proxies of Atari obscure.

In this work we use Pok\'emon Red as a controlled test bed for one
specific question: \textbf{how does the choice of observation
representation affect what a long-horizon RL agent can learn?}  We
compare three treatments under matched architecture, reward, and
compute:
%
\begin{itemize}
\item \textbf{Pixel} --- the standard $80 \times 72$ grayscale screen
  observation, encoded with a Nature-DQN-style
  CNN~\citep{mnih2015atari}.
\item \textbf{Symbolic} --- a structured vector read directly from the
  emulator RAM, containing the local tile map, the 18 pre-registered
  event flags, party stats, and player coordinates, encoded with an
  MLP.
\item \textbf{Hybrid} --- both streams concatenated at the recurrent
  layer.
\end{itemize}
%
All three feed into a single LSTM and a PPO actor-critic
head~\citep{schulman2017ppo,hausknecht2015drqn}, with the same
event-flag-based reward function.  Differences in performance are
attributable to the observation representation alone.

\paragraph{Contributions.}
The contributions of this preliminary study, intended to seed a longer
journal-length follow-up, are:
\begin{enumerate}
\item \textbf{An open benchmark.}  We release \textsc{PokeGym}, a
  Gymnasium-compatible~\citep{gymnasium} environment for Pok\'emon Red
  with three first-class observation treatments, an event-flag reward
  calculator, and a curriculum-aware save-state system.  The
  environment runs headlessly at $\sim$\todopilot{N}~steps/s/CPU on a
  consumer laptop.
\item \textbf{A pre-registered analysis plan.}  Hypotheses, primary
  metric (Boulder Badge win-rate), seed counts, statistical tests
  (\texttt{rliable} stratified bootstrap with 95\% CIs), and stopping
  rules were committed to version control before the main experiments
  ran.\footnote{See \texttt{paper/analysis\_plan.md} in the open
    repository.}
\item \textbf{Empirical comparison of three observation treatments at
  10M~environment steps each, three seeds per treatment.}
  \todopilot{Replace this bullet once results are in: 1-2 sentences on
    the dominant treatment and the magnitude of the gap, with the
    bootstrap CI.}
\end{enumerate}

\paragraph{Scope.}
The pilot results reported here use 10M environment steps per seed and
three seeds per treatment.  We treat them as preliminary evidence
rather than a final ranking---both the seed count and the training
budget will be substantially expanded in subsequent journal-length
work, as committed in the pre-registered plan.

The remainder of the paper is organised as follows.
\Cref{sec:related-work} situates the work in the long-horizon-RL,
representation-learning, and Pok\'emon-RL literatures.
\Cref{sec:environment} describes \textsc{PokeGym} and the three
observation treatments in detail.  \Cref{sec:methods} fixes the
training pipeline, statistical methodology, and pre-registered
deviations.  \Cref{sec:results} presents pilot learning curves,
aggregate metrics, and milestone progression.  \Cref{sec:discussion}
discusses what the pilot can and cannot conclude, and outlines the
follow-up experiments planned for the journal version.
